\documentclass[a4paper,11pt]{article}
        \usepackage[top=15mm,bottom=18mm,left=16mm,right=16mm]{geometry}
        \usepackage{fontspec}
        \usepackage{xeCJK}
        \setmainfont{Times New Roman}
        \setCJKmainfont{Yu Gothic}
        \setmonofont{Consolas}
        \setCJKmonofont{Yu Gothic}
        \usepackage{booktabs}
        \usepackage{longtable}
        \usepackage{tabularx}
        \usepackage{array}
        \usepackage{enumitem}
        \usepackage{hyperref}
        \usepackage{listings}
        \usepackage{xcolor}
        \hypersetup{
          colorlinks=true,
          linkcolor=blue,
          urlcolor=blue,
          pdftitle={上位MPC 距離メッシュ生成},
          pdfauthor={Codex}
        }
        \lstset{
          basicstyle=\ttfamily\small,
          breaklines=true,
          columns=fullflexible,
          keepspaces=true,
          frame=single,
          framerule=0.2pt,
          rulecolor=\color{black},
          showstringspaces=false
        }
        \setlength{\parskip}{0.42em}
        \setlength{\parindent}{1em}
        \renewcommand{\arraystretch}{1.15}
        \title{14. 上位MPC 距離メッシュ生成}
        \author{solar\_ws0129-main}
        \date{2026-07-29}
        \begin{document}
        \maketitle

        \section*{対象}
        \begin{tabularx}{\linewidth}{>{\raggedright\arraybackslash}p{0.18\linewidth} X}
        \toprule
        項目 & 内容 \\
        \midrule
        ファイル & \path|mpc_solarcar/upper_horizon.py| \\
        種別 & Python \\
        区分 & planner helper \\
        役割 & 固定または適応距離メッシュを作り、現在地点から先の control point と segment を決める。 \\
        起動文脈 & distance-domain upper planner の最初の一歩。 \\
        \bottomrule
        \end{tabularx}

        \section*{このファイルを読む理由}
        固定または適応距離メッシュを作り、現在地点から先の control point と segment を決める。

        \begin{itemize}[leftmargin=1.6em]
        \item build\_upper\_distance\_horizon が中心。
\item plan\_segment\_index が現在位置から有効速度区間を引く。
        \end{itemize}

        \section*{呼び出し関係}
        \subsection*{どこから呼ばれるか}
        \begin{itemize}[leftmargin=1.6em]
        \item \path|mpc_node.py|
\item \path|scripts/solar_sim.py|
        \end{itemize}

        \subsection*{次に読むと理解が進むファイル}
        \begin{itemize}[leftmargin=1.6em]
        \item \path|mpc_solarcar/upper_policy.py|
        \end{itemize}

        \subsection*{同一リポジトリ内の主要依存}
        \begin{itemize}[leftmargin=1.6em]
        \item 同一リポジトリ内の明示 import は少ない。
        \end{itemize}

        \section*{ソース冒頭の把握}
        モジュール docstring または先頭意図の要約:

        モジュール docstring は明示されていない。

        \subsection*{代表コード断片}
        \begin{lstlisting}[language=Python]
@dataclass
class UpperDistanceHorizon:
    ds_seq_km: np.ndarray
    seg_s_km: np.ndarray
    ctrl_s_km: np.ndarray

    @property
    def total_km(self) -> float:
        return float(np.sum(self.ds_seq_km)) if len(self.ds_seq_km) else 0.0


def _weighted_extra(count: int, growth: float) -> np.ndarray:
    if count <= 1:
        return np.ones(max(1, count), dtype=float)
    growth = max(1.0, float(growth))
    if abs(growth - 1.0) <= 1.0e-9:
        return np.ones(count, dtype=float)
    return np.power(growth, np.arange(count, dtype=float))
\end{lstlisting}


        \section*{主要構造の読み取り}
        主要クラスは UpperDistanceHorizon。 主要関数は total\_km, build\_upper\_distance\_horizon, plan\_segment\_index。

        \subsection*{主要クラス・関数}
            \begin{longtable}{p{0.07\linewidth} p{0.16\linewidth} p{0.25\linewidth} p{0.40\linewidth}}
            \toprule
            行 & 種別 & 名前 & 補足 \\
            \midrule
            14 & class & \path|UpperDistanceHorizon| &  \\
20 & function & \path|total_km| & args: self \\
24 & function & \path|_weighted_extra| & args: count, growth \\
33 & function & \path|_adaptive_ds_sequence| & args: target\_km \\
82 & function & \path|build_upper_distance_horizon| &  \\
136 & function & \path|plan_segment_index| & args: plan\_segments, s\_km \\
            \bottomrule
            \end{longtable}

        \subsection*{ファイルを上から読んだときの定義順}
            \begin{longtable}{p{0.07\linewidth} p{0.84\linewidth}}
            \toprule
            行 & 内容 \\
            \midrule
            10 & MIN\_DISTANCE\_STEP\_KM に 1e-06 の結果を代入する。 \\
14 & クラス UpperDistanceHorizon を定義する。 \\
24 & 関数 \_weighted\_extra を定義する。 \\
33 & 関数 \_adaptive\_ds\_sequence を定義する。 \\
82 & 関数 build\_upper\_distance\_horizon を定義する。 \\
136 & 関数 plan\_segment\_index を定義する。 \\
            \bottomrule
            \end{longtable}

        \subsection*{import 群の詳細}
            \begin{longtable}{p{0.07\linewidth} p{0.33\linewidth} p{0.50\linewidth}}
            \toprule
            行 & import 文 & このファイルで読む理由 \\
            \midrule
            1 & \path|from __future__ import annotations| & 型ヒントの遅延評価を有効にし、自己参照型や forward reference を安全に扱うため。 このファイル内での使用位置は少ないか、間接利用である。 \\
3 & \path|import math| & clip、sqrt、sin/cos、有限判定など数値ロジックに使うため。 このファイル内での主な使用位置は L48, L112。 \\
4 & \path|from dataclasses import dataclass| & dataclasses から dataclass を読み込み、このファイルの処理を組み立てるため。 このファイル内での主な使用位置は L13。 \\
5 & \path|from typing import Iterable, List| & typing から Iterable, List を読み込み、このファイルの処理を組み立てるため。 このファイル内での主な使用位置は L136, L137。 \\
7 & \path|import numpy as np| & 数値配列、clip、補間、統計計算を行うため。 このファイル内での主な使用位置は L15, L16, L17, L21, L24, L26, L29, L30, ...。 \\
            \bottomrule
            \end{longtable}

        \section*{関数・クラスを上から順に解説}
        \subsection*{L14 クラス \path|UpperDistanceHorizon|}
            \begin{itemize}[leftmargin=1.6em]
              \item 定義: \path|UpperDistanceHorizon(bases=なし)|
              \item docstring: 明示されていない。
              \item このブロックが直接呼ぶ主な関数・メソッド: float, len, sum
              \item 戻り値の要点: 戻り値が無いか、return 文が明示されていない。
            \end{itemize}
            \begin{enumerate}[leftmargin=1.8em]
            \item ds\_seq\_km に  を代入する。
\item seg\_s\_km に  を代入する。
\item ctrl\_s\_km に  を代入する。
\item 関数 total\_km を定義する。
            \end{enumerate}

\subsection*{L24 関数 \path|_weighted_extra|}
            \begin{itemize}[leftmargin=1.6em]
              \item 定義: \path|_weighted_extra(count, growth)|
              \item docstring: 明示されていない。
              \item このブロックが直接呼ぶ主な関数・メソッド: abs, arange, float, max, ones, power
              \item 戻り値の要点: np.power(growth, np.arange(count, dtype=float)) / np.ones(max(1, count), dtype=float) / np.ones(count, dtype=float)
            \end{itemize}
            \begin{enumerate}[leftmargin=1.8em]
            \item 条件 count <= 1 を判定し、真なら内部処理を行う。
\item growth に max(1.0, float(growth)) の結果を代入する。
\item 条件 abs(growth - 1.0) <= 1e-09 を判定し、真なら内部処理を行う。
\item np.power(growth, np.arange(count, dtype=float)) を返す。
            \end{enumerate}

\subsection*{L33 関数 \path|_adaptive_ds_sequence|}
            \begin{itemize}[leftmargin=1.6em]
              \item 定義: \path|_adaptive_ds_sequence(target_km, max_steps, min_ds_km, max_ds_km, growth)|
              \item docstring: 明示されていない。
              \item このブロックが直接呼ぶ主な関数・メソッド: \_weighted\_extra, any, array, ceil, copy, enumerate, flatnonzero, float, full, int, max, maximum
              \item 戻り値の要点: ds / np.array([target\_km], dtype=float) / np.array([target\_km], dtype=float) / np.full(step\_count, target\_km / step\_count, dtype=float)
            \end{itemize}
            \begin{enumerate}[leftmargin=1.8em]
            \item target\_km に max(MIN\_DISTANCE\_STEP\_KM, float(target\_km)) の結果を代入する。
\item min\_ds\_km に max(MIN\_DISTANCE\_STEP\_KM, float(min\_ds\_km)) の結果を代入する。
\item 条件 max\_steps <= 1 を判定し、真なら内部処理を行う。
\item 条件 target\_km <= min\_ds\_km を判定し、真なら内部処理を行う。
\item step\_count に int(min(max\_steps, max(1, math.ceil(target\_km / min\_ds\_km)))) の結果を代入する。
\item base に np.full(step\_count, min\_ds\_km, dtype=float) の結果を代入する。
\item base\_sum に float(base.sum()) の結果を代入する。
\item 条件 target\_km <= base\_sum + 1e-09 を判定し、真なら内部処理を行う。
\item weights に \_weighted\_extra(step\_count, growth) の結果を代入する。
\item weights に weights / max(float(weights.sum()), 1e-09) の結果を代入する。
            \end{enumerate}

\subsection*{L82 関数 \path|build_upper_distance_horizon|}
            \begin{itemize}[leftmargin=1.6em]
              \item 定義: \path|build_upper_distance_horizon(mode, s0_km, race_km, ds_km, horizon_km, max_steps, ctrl_km, adaptive_min_ds_km, adaptive_max_ds_km, adaptive_growth)|
              \item docstring: 明示されていない。
              \item このブロックが直接呼ぶ主な関数・メソッド: UpperDistanceHorizon, \_adaptive\_ds\_sequence, append, arange, array, ceil, concatenate, cumsum, float, full, int, len
              \item 戻り値の要点: UpperDistanceHorizon(ds\_seq\_km=ds\_seq, seg\_s\_km=seg\_s, ctrl\_s\_km=ctrl\_s)
            \end{itemize}
            \begin{enumerate}[leftmargin=1.8em]
            \item ds\_km に max(MIN\_DISTANCE\_STEP\_KM, float(ds\_km)) の結果を代入する。
\item horizon\_km に max(ds\_km, float(horizon\_km)) の結果を代入する。
\item remaining\_km に max(0.0, float(race\_km) - float(s0\_km)) の結果を代入する。
\item max\_steps に max(1, int(max\_steps)) の結果を代入する。
\item mode に str(mode or 'fixed').strip().lower() の結果を代入する。
\item 条件 mode in \{'adaptive\_full\_race', 'remaining\_race', 'full\_race'\} を判定し、真なら内部処理を行う。
\item seg\_s に np.concatenate(([0.0], np.cumsum(ds\_seq[:-1], dtype=float))) の結果を代入する。
\item total\_km に float(np.sum(ds\_seq)) の結果を代入する。
\item control\_end\_km に float(seg\_s[-1]) if len(seg\_s) else 0.0 の結果を代入する。
\item ctrl\_step\_km に float(ctrl\_km) if ctrl\_km and ctrl\_km > 0.0 else float(ds\_seq[0]) の結果を代入する。
            \end{enumerate}

\subsection*{L136 関数 \path|plan_segment_index|}
            \begin{itemize}[leftmargin=1.6em]
              \item 定義: \path|plan_segment_index(plan_segments, s_km)|
              \item docstring: 明示されていない。
              \item このブロックが直接呼ぶ主な関数・メソッド: enumerate, float, get, len, list
              \item 戻り値の要点: len(segments) - 1 / -1 / idx
            \end{itemize}
            \begin{enumerate}[leftmargin=1.8em]
            \item segments に list(plan\_segments) を代入する。
\item 条件 not segments を判定し、真なら内部処理を行う。
\item s\_km に float(s\_km) の結果を代入する。
\item enumerate(segments) を順に走査し、各要素を (idx, seg) に入れて処理する。
\item len(segments) - 1 を返す。
            \end{enumerate}











        \section*{処理の流れ}
        \begin{enumerate}[leftmargin=1.7em]
        \item ソース中の主要関数を通じて処理を進める。
        \end{enumerate}

        \section*{このファイルの位置づけ}
        この資料群では、\path|mpc_solarcar/upper_horizon.py| を
        \textbf{planner helper}の中核として扱う。
        したがって、ここで扱う処理は単独で完結せず、
        前段の profile / launch / telemetry と後段の planner / logger / report へ繋がる。

        \end{document}
